\section{实验条件}
% 11
\begin{frame}[t]{不同医学 + AI 课题的实验条件：一张定位图}
\vspace{.05cm}
\centering
\begin{tikzpicture}[x=1.1cm,y=.38cm,
  pt/.style={circle,fill=Muted,inner sep=2.2pt},
  me/.style={circle,fill=DarkBlue,inner sep=2.6pt},
  lb/.style={anchor=west,align=left,font=\scriptsize,inner sep=2pt}]
 \draw[-{Latex[length=2mm]},draw=Muted,line width=.6pt] (0,0) -- (10.6,0) node[below left,font=\scriptsize,text=Muted]{数据获取与合规门槛};
 \draw[-{Latex[length=2mm]},draw=Muted,line width=.6pt] (0,0) -- (0,10.6) node[anchor=south west,font=\scriptsize,text=Muted,xshift=-2pt]{算力与存储需求};
 \draw[dashed,draw=Line] (5,0) -- (5,10); \draw[dashed,draw=Line] (0,5) -- (10,5);
 \node[pt] (sig) at (2.6,1.2) {}; \node[lb] at (sig.east) {\textbf{\faHeartbeat\ 生理信号}\\\color{Muted}长时序、标签稀疏；时序网络};
 \node[me] (o1) at (1.2,3.4) {}; \node[lb] (o1l) at (o1.east) {\textbf{\color{DarkBlue}\faDna\ 组学 / 虚拟细胞 · 起步}\\\color{Muted}公开 Perturb-seq；CPU 基线 + 单卡微调};
 \node[pt] (txt) at (7.0,3.2) {}; \node[lb] at (txt.east) {\textbf{\faIcon{file-medical-alt}\ 临床文本}\\\color{Muted}中文脱敏难；LLM LoRA 微调单卡};
 \node[pt] (img) at (4.2,7.0) {}; \node[lb] at (img.east) {\textbf{\faXRay\ 医学影像}\\\color{Muted}需医生标注；3D 体数据显存高};
 \node[me] (o2) at (8.6,6.0) {}; \node[lb] at (o2.east) {\textbf{\color{DarkBlue}\faDna\ 组学 / 虚拟细胞 · 合作}\\\color{Muted}院方样本；伦理与人遗审批；院内 GPU};
 \node[pt] (pat) at (8.3,8.8) {}; \node[lb] at (pat.east) {\textbf{\faMicroscope\ 数字病理}\\\color{Muted}GB 级切片；存储与多卡};
 \draw[-{Latex[length=2mm]},dashed,draw=DarkBlue,line width=.8pt] ([xshift=2pt]o1l.east) -- (o2);
\end{tikzpicture}
\remark{起步门槛最低是本次方向的优势；进入合作阶段后，门槛上升到与影像、文本课题相当。}
\source{经验性定位，用于讨论合作条件；位置与显存需按课题、模型与数据规模核对，不作为文献结论。}
\end{frame}
